% GENERATED FILE. Edit canonical lessons, then run npm run generate:books.
% canonical-source-hash: fnv1a64:0a5fe5023bf46961
% canonical-lessons: JA-C25-tomodachi, JA-C25-otona, JA-C25-otoko, JA-C25-onna, JA-C25-sensei

\chapter{Friend, Adult, Man, Woman, Teacher}
\label{ch:ja-friend-adult-man-woman-teacher}
\hlchaptermodality{25}

\begin{chapteropening}
\textbf{By the end of this chapter:} \emph{I can say a friend, an adult, a man, a woman and a teacher.}

Complete the last lesson of chapter 25: i can say a friend, an adult, a man, a woman and a teacher.
\end{chapteropening}

\section[tomodachi]{\ja{ともだち} (tomodachi) — a friend}

\label{lesson:JA-C25-tomodachi}



\begin{quote}
Before the new one: say the Japanese for a son, then the Japanese for a daughter.
\end{quote}

\subsection*{You'll want to know: \ja{ともだち}}
\textbf{\ja{ともだち}} — \emph{tomodachi} — \textquotedblleft{}a friend\textquotedblright{}.

\begin{grammarlens}[title={What you've built}]
One more word for this chapter.
\end{grammarlens}

\subsection*{Your turn}
\begin{itemize}
\raggedright
  \item \emph{Say it:} \emph{tomodachi}
  \item \emph{Say it:} \emph{tomodachi}, once more
  \item \emph{Say it:} say \emph{tomodachi}
  \item \emph{Recall:} say the Japanese for a son, then the Japanese for a daughter, then say \emph{tomodachi} again
\end{itemize}

\begin{tcolorbox}[breakable,colback=teal!4,colframe=teal!35!black,title={Before you move on}]
What does \ja{ともだち} mean? (\textquotedblleft{}A friend\textquotedblright{}.) Say it once more.
\end{tcolorbox}
\section[otona]{\ja{おとな} (otona) — an adult}

\label{lesson:JA-C25-otona}



\begin{quote}
Before the new one: say the Japanese for a daughter, then the Japanese for a friend.
\end{quote}

\subsection*{You'll want to know: \ja{おとな}}
\textbf{\ja{おとな}} — \emph{otona} — \textquotedblleft{}an adult\textquotedblright{}.

\begin{grammarlens}[title={What you've built}]
One more word for this chapter.
\end{grammarlens}

\subsection*{Your turn}
\begin{itemize}
\raggedright
  \item \emph{Say it:} \emph{otona}
  \item \emph{Say it:} \emph{otona}, once more
  \item \emph{Say it:} say \emph{otona}
  \item \emph{Recall:} say the Japanese for a daughter, then the Japanese for a friend, then say \emph{otona} again
\end{itemize}

\begin{tcolorbox}[breakable,colback=teal!4,colframe=teal!35!black,title={Before you move on}]
What does \ja{おとな} mean? (\textquotedblleft{}An adult\textquotedblright{}.) Say it once more.
\end{tcolorbox}
\section[otoko]{\ja{おとこ} (otoko) — a man}

\label{lesson:JA-C25-otoko}



\begin{quote}
Before the new one: say the Japanese for a friend, then the Japanese for an adult.
\end{quote}

\subsection*{You'll want to know: \ja{おとこ}}
\textbf{\ja{おとこ}} — \emph{otoko} — \textquotedblleft{}a man\textquotedblright{}.

\begin{grammarlens}[title={What you've built}]
One more word for this chapter.
\end{grammarlens}

\subsection*{Your turn}
\begin{itemize}
\raggedright
  \item \emph{Say it:} \emph{otoko}
  \item \emph{Say it:} \emph{otoko}, once more
  \item \emph{Say it:} say \emph{otoko}
  \item \emph{Recall:} say the Japanese for a friend, then the Japanese for an adult, then say \emph{otoko} again
\end{itemize}

\begin{tcolorbox}[breakable,colback=teal!4,colframe=teal!35!black,title={Before you move on}]
What does \ja{おとこ} mean? (\textquotedblleft{}A man\textquotedblright{}.) Say it once more.
\end{tcolorbox}
\section[onna]{\ja{おんな} (onna) — a woman}

\label{lesson:JA-C25-onna}



\begin{quote}
Before the new one: say the Japanese for an adult, then the Japanese for a man.
\end{quote}

\subsection*{You'll want to know: \ja{おんな}}
\textbf{\ja{おんな}} — \emph{onna} — \textquotedblleft{}a woman\textquotedblright{}.

\begin{grammarlens}[title={What you've built}]
One more word for this chapter.
\end{grammarlens}

\subsection*{Your turn}
\begin{itemize}
\raggedright
  \item \emph{Say it:} \emph{onna}
  \item \emph{Say it:} \emph{onna}, once more
  \item \emph{Say it:} say \emph{onna}
  \item \emph{Recall:} say the Japanese for an adult, then the Japanese for a man, then say \emph{onna} again
\end{itemize}

\begin{tcolorbox}[breakable,colback=teal!4,colframe=teal!35!black,title={Before you move on}]
What does \ja{おんな} mean? (\textquotedblleft{}A woman\textquotedblright{}.) Say it once more.
\end{tcolorbox}
\section[sensei]{\ja{せんせい} (sensei) — a teacher}

\label{lesson:JA-C25-sensei}



\begin{quote}
Before the new one: say the Japanese for a man, then the Japanese for a woman.
\end{quote}

\subsection*{You'll want to know: \ja{せんせい}}
\textbf{\ja{せんせい}} — \emph{sensei} — \textquotedblleft{}a teacher\textquotedblright{}.

\begin{grammarlens}[title={What you've built}]
One more word for this chapter.
\end{grammarlens}

\subsection*{Your turn}
\begin{itemize}
\raggedright
  \item \emph{Say it:} \emph{sensei}
  \item \emph{Say it:} \emph{sensei}, once more
  \item \emph{Say it:} say \emph{sensei}
  \item \emph{Recall:} say the Japanese for a man, then the Japanese for a woman, then say \emph{sensei} again
\end{itemize}

\begin{tcolorbox}[breakable,colback=teal!4,colframe=teal!35!black,title={Before you move on}]
What does \ja{せんせい} mean? (\textquotedblleft{}A teacher\textquotedblright{}.) Say it once more.
\end{tcolorbox}
